% 来源: 19c9fb8e-8112-85ce-8000-0000310bad72_2.2_传输媒体.pdf
% 提取时间: 2026-02-28T22:51:31+08:00

2.2 传输媒体
有线传输媒体 or 导向媒体
 双绞线（Twisted Pair）
 同轴电缆（Coaxial Cable）
 光纤（Optical Fiber）
 常用的线缆高速连接模块与接口
  直接连接铜缆 （Direct-Attach Copper, DAC）
  有源光缆（AOC, Active Optical Cable）
无线媒体 or 非导向媒体
 按照频率特性的信号划分

传输媒体是发送器和接收器之间的物理通道， 首先这里给一个具像化的表示，可以看得更清楚一
些




通过电磁波频谱分布特征来看一下对应的应用。电磁波频谱是电磁波按频率或波长划分的范围，广
泛应用于通信、广播、雷达、医疗等领域。电磁波频谱从极低频（ELF）到极高频（EHF）以及更
高频段，每个频段都有其独特的特性和应用。
以几个主要频段为例。
 ● 在极低频（3 Hz ~ 30 Hz）和超低频（30 Hz ~ 300 Hz）具有极强的穿透能力，适合水下和
   地下通信，常用于潜艇通信和地球物理探测。也是话音和乐器所在的频段。
 ● 低频（LF，30 kHz ~ 300 kHz）和中频（MF，300 kHz ~ 3 MHz）主要依靠地面波传播，适
   合长距离广播和导航系统，如AM广播电台。
 ● 高频（HF，3 MHz ~ 30 MHz）利用电离层反射实现远距离通信，广泛应用于短波广播和航空
   通信。
 ● 甚高频（VHF，30 MHz ~ 300 MHz）和特高频（UHF，300 MHz ~ 3 GHz）则以直线传播
   为主，适合电视广播、FM广播、移动通信（如4G）和无线局域网（Wi-Fi）。
 ● 超高频（SHF，3 GHz ~ 30 GHz）和极高频（EHF，30 GHz ~ 300 GHz）具有极高的带宽
   和传输速率，但传播距离较短，主要用于卫星通信、5G网络和雷达系统。
 ● 红外线、可见光、紫外线等光波频段也在自由空间光通信、激光通信和医疗领域有重要应
   用。
上图同样给出了物理媒体的工作区间。根据其特性和应用场景，分为有线传输媒体（双绞线、同轴
电缆、光纤）和无线传输媒体（视距传播、地波、天波）两大类
有线传输媒体 or 导向媒体
双绞线（Twisted Pair）
双绞线是一种常用的有线传输媒体，由两根绝缘的铜导线以螺旋形式绞合在一起组成。这种设计可
以减少电磁干扰（EMI）和串扰，提高信号传输质量。
线缆间串扰：信号通道之间的耦合现象。 双绞线的扭绞结构是为了减少相邻导线之间的串扰和消
除外界干扰
双绞线广泛应用于局域网（LAN）、电话系统和其他短距离通信场景。
结构
 ● 导线：两根绝缘的铜导线。
 ● 绞合方式：两根导线以一定绞距螺旋绞合，以减少外部电磁干扰和线对间的串扰。
 ● 外护套：外层包裹绝缘材料（如PVC）
                    ，用于保护内部导线并提供机械强度。




双绞线
                           网线
双绞线根据是否屏蔽和性能等级分为多种类型：
(1) 非屏蔽双绞线（UTP, Unshielded Twisted Pair）
  ● 无屏蔽层，成本低，易于安装。
  ● 抗干扰能力较弱，适合低噪声环境。
  ● 应用：以太网（如Cat5e、Cat6）    、电话线。
(2) 屏蔽双绞线（STP, Shielded Twisted Pair）
  ● 在绞合线对外增加金属屏蔽层（如铝箔或编织网）              ，抗干扰能力强。
  ● 成本较高，安装复杂。
  ● 应用：工业环境、高噪声场景。
双绞线作为一种常用的传输媒体，具有成本低、易于安装和维护的优点，其绞合设计能有效减少电
磁干扰和串扰，适合短距离通信和局域网部署。然而，双绞线的传输距离有限（通常不超过100
米），带宽较低，且非屏蔽双绞线（UTP）在高噪声环境中的抗干扰能力较弱。尽管如此，双绞线
凭借其经济性和广泛适用性，仍然是许多通信场景中的重要选择。
同轴电缆（Coaxial Cable）
同轴电缆是一种常用的有线传输媒体，由内导体、绝缘层、屏蔽层和外护套组成。其结构设计使其
具有较高的带宽和抗干扰能力，广泛应用于电视广播、网络通信和视频监控等领域。
1. 同轴电缆的结构
  ● 内导体：中心为铜或铝制成的导线，用于传输信号。

  ● 绝缘层：包裹内导体的绝缘材料（如聚乙烯），用于隔离内导体和屏蔽层。
  ● 屏蔽层：金属编织网或铝箔，用于屏蔽外部电磁干扰。
  ● 外护套：最外层的绝缘材料（如PVC）
                     ，用于保护电缆并提供机械强度。




2. 同轴电缆的分类
同轴电缆根据阻抗和用途分为多种类型：
(1) 射频同轴电缆
  ● 阻抗：50Ω或75Ω。

  ● 应用：电视广播、无线电通信、卫星通信。

(2) 视频同轴电缆
  ● 阻抗：75Ω。
  ● 应用：闭路电视（CCTV）
                、视频监控。
(3) 网络同轴电缆
  ● 阻抗：50Ω。
  ● 应用：早期以太网（如10BASE2、10BASE5）
                             。
光纤（Optical Fiber）
光纤是一种利用光信号传输数据的介质，由高纯度的玻璃或塑料制成。由于其高带宽、低损耗和抗
干扰能力强，光纤已成为现代通信网络的核心组成部分，广泛应用于数据中心、电信网络和互联网
骨干网等领域。高锟1966年的论文分析了玻璃纤维损耗大的主要原因，并大胆预言只要能设法降
低玻璃纤维的杂质，就有可能使光纤的损耗从每公里1000dB降低到20dB，从而用于通信。2009
年高锟获得诺贝尔物理学奖。1970年美国康宁公司设计和演示了低损耗18dB/km的光纤。
1. 光纤的结构
光纤通常由以下几层组成：
  ● 纤芯（Core）：中心部分，光信号在此传输。纤芯的直径通常为8～62.5微米。
  ● 包层（Cladding）：包裹纤芯的材料，折射率低于纤芯，用于将光信号限制在纤芯内。
  ● 涂覆层（Coating）：保护光纤的外层，通常由塑料制成，防止机械损伤和环境影响。




2. 光纤的类型
根据传输模式和材料特性，光纤可以分为以下几类：
(1) 单模光纤（Single-Mode Fiber, SMF）单模光纤只允许一种模式的光在其中传播，即基模。
这种模式的光沿着光纤的中心轴线传播，没有模式色散，因此可以实现长距离、高速率的数据传
输，适用于大容量、长距离的通信系统。
  ● 纤芯直径小（通常为8～10微米）        ，只允许一种模式的光信号传输。
  ● 传输距离远（可达数十公里到上百公里）           ，带宽极高。
  ● 成本较高，适合长距离和高带宽应用。

  ● 应用：电信骨干网。数据中心互联（DCI）          。长距离通信（如海底光缆）。
(2) 多模光纤（Multi-Mode Fiber, MMF）多模光纤允许多种模式的光同时在纤芯中传播。不同模
式的光在光纤中传播的路径和速度略有不同，这就会导致模式色散，使得光信号在传输过程中发生
展宽，限制了传输距离和传输速率。多模光纤常用于短距离、低速率的通信系统
  ● 纤芯直径较大（通常为50～62.5微米）        ，允许多种模式的光信号传输。
  ● 传输距离较短（通常为几百米到2公里）          ，带宽较低。
  ● 成本较低，适合短距离和低成本应用。
  ● 应用：局域网（LAN）   。数据中心内部连接。短距离通信（如校园网）。
3.   光纤的工作区




目前光纤通信的主要波长工作区包括：
  ● 850nm： 用于多模光纤，适合短距离通信。
  ● 1310nm： 用于单模光纤，适合中距离通信。

  ● 1550nm： 用于单模光纤，适合长距离通信，支持波分复用技术。

这些波长区域的选择取决于传输距离、带宽需求和成本等因素。随着波分复用技术的发展，光纤的
传输容量和效率将进一步提升，为未来高速通信网络提供强大支持。
4. 光纤的传输原理
光纤利用全反射原理传输光信号：当光从光密介质（折射率较大）射向光疏介质（折射率较小）
时，且入射角大于临界角，就会发生全反射现象。在光纤中，纤芯的折射率高于包层的折射率，光
在纤芯中传播时，遇到纤芯与包层的界面，只要入射角满足条件，光就会在纤芯内不断地发生全反
射，从而沿着光纤向前传播，就像光在光纤内部 “跳跃” 前进一样。
 1. 光信号从纤芯的一端入射。
 2. 由于纤芯的折射率高于包层，光信号在纤芯内不断反射，向前传播。
 3. 光信号最终到达光纤的另一端，完成数据传输。

4. 光纤的优点
 ●  高带宽：支持Tbps级别的传输速率，远超铜缆。
  ● 低损耗：光信号在光纤中传输时衰减极小，适合长距离传输。
  ● 抗干扰能力强：不受电磁干扰，适合高噪声环境。
  ● 安全性高：光信号难以被窃听，数据传输更安全。

  ● 重量轻、体积小：光纤比铜缆更轻便，适合高密度部署。

光纤是否可以弯曲？
光纤在设计和应用中可以弯曲，但其弯曲性能受到严格的物理限制。光纤的弯曲半径必须控制在最
小弯曲半径以内，以避免光信号从纤芯泄漏到包层，导致宏弯损耗或微弯损耗。单模光纤的最小弯
曲半径通常为10倍光纤直径（约1.25毫米），而多模光纤则为20倍光纤直径（约2.5毫米）。过度弯
曲不仅会增加信号衰减，还可能导致光纤断裂，影响传输性能。为了解决这一问题，近年来开发了
弯曲不敏感光纤（Bend-Insensitive Fiber），通过优化纤芯和包层的折射率分布，显著降低了弯曲
损耗，使其在更小的弯曲半径（通常≤5毫米）下仍能保持低损耗和高性能。这种光纤特别适用于
高密度布线、复杂安装环境以及光纤到户（FTTH）等场景。在实际应用中，安装光纤时需严格遵
守最小弯曲半径的要求，并使用保护套管或光纤管理设备来防止外力损伤。定期检查光纤的弯曲状
态也是确保其长期稳定运行的重要措施。总之，光纤的弯曲性能是其设计和部署中需要重点关注的
技术参数，合理的弯曲管理能够有效保障光纤通信系统的高效性和可靠性。
常用的光源有哪些？
常用的光纤通信光源主要包括激光二极管（LD, Laser Diode）和发光二极管（LED, Light
Emitting Diode），它们发射的光波长通常位于光纤的低损耗窗口，如850nm、1310nm和
1550nm。
  ● 激光二极管因其高功率、窄光谱和高速调制特性，成为长距离、高带宽光纤通信的首选光源，
    主要分为分布反馈激光器（DFB Laser）和垂直腔面发射激光器（VCSEL）。DFB激光器发射
    的光波长通常为1310nm或1550nm，具有极高的单色性和稳定性，适用于单模光纤的长距离传
    输，尤其是在电信骨干网和数据中心互联（DCI）中广泛应用；而VCSEL则主要发射850nm波
    长的光，因其低成本、低功耗和高调制速率，广泛用于多模光纤的短距离通信，如数据中心内
    部互联和高性能计算集群。
  ● 发光二极管（LED）虽然成本低且易于制造，但其光谱较宽、调制速率较低，通常发射
    850nm或1310nm波长的光，适用于低速率、短距离的光纤通信系统，如早期的局域网（LAN）
    和光纤到户（FTTH）应用。
此外，随着技术的进步，量子点激光器和外调制激光器等新型光源也逐渐崭露头角，它们通过发射
更稳定的光波长（如1550nm）和提升调制速率，为下一代高速光纤通信系统提供了强有力的支
持。
常用的线缆高速连接模块与接口
SFP（Small Form-factor Pluggable） 是一种小型可插拔的光模块或电模块，广泛用于网络设
备（如交换机、路由器和网卡）中，用于实现光纤或铜缆的连接。
SFP模块的后续版本如下表，这些模块和接口标准推动了数据中心和高性能计算网络的发展，满足
了日益增长的带宽需求。
 类型     速率      通道数       特点         应用
 SFP+ 10Gbps 单通道          体积小，功耗低 10G以太网网络
 SFP28 25Gbps 单通道         支持25G以太网， 25G以太网网络
                          兼容SFP+
 QSFP+ 40Gbps 4通道         支持40G以太网和 40G以太网网络，HPC集群
                （10Gbps/通 InfiniBand
                道）
 QSFP28 100Gbps 4通道       支持100G以太网， 100G以太网网络，数据中心
                （25Gbps/通 兼容QSFP+    骨干网络
                道）
 QSFP56 200Gbp 4通道        支持200G以太网， 200G以太网网络，超大规模
        s       （50Gbps/通 兼容QSFP28   数据中心
                道）
 QSFP- 400Gbp 8通道         支持400G以太网， 400G以太网网络，下一代数
 DD     s       （50Gbps/通 双排引脚设计     据中心互联
                道）
注：SFP28 的单通道速率为28 Gbps，但实际用于25 Gigabit Ethernet（25GbE）时，数据传输
速率为25 Gbps。额外的3 Gbps 用于编码开销（如FEC，前向纠错），以确保数据传输的可靠
性。QSFP-DD 中，“DD” 表示 Double Density（双密度），指的是该光模块的接口设计支持双
排引脚，从而实现了更高的端口密度和传输带宽。QSFP-DD 支持 8 通道，每通道速率可达 50
Gbps（NRZ 调制）或 100 Gbps（PAM4 调制）。总带宽可达 400 Gbps（8x50 Gbps） 或 800
Gbps（8x100 Gbps）。
直接连接铜缆 （Direct-Attach Copper, DAC）
一种用于小型可插拔（SFP）以太网的标准电缆类型，最初定义为SFP+直接连接铜缆
（10GSFP+Cu），它通过有源或无源双轴电缆组件提供10Gbps以太网，并直接连接到SFP+接口。
有源双轴电缆在SFP+接口中包含有源电子元件以改善信号质量；而无源双轴电缆主要是一条直通
的“电线”，几乎不包含任何元件。通常，长度小于7米的双轴电缆是无源的，而长度超过7米的则
是有源的，但这可能因供应商而异。SFP+直接连接铜缆（DAC）因其低延迟和低成本，成为
10Gbps以太网传输距离达10米的流行选择。




结构：
  ● 内部导线：多对双绞线或平行铜线。

  ● 连接器：集成高速连接器（如QSFP+、QSFP28、QSFP-DD）       。
  ● 屏蔽层：通常采用金属屏蔽层减少电磁干扰。

  ● 外护套：柔性外护套，便于布线和安装。

特点：
  ● 无源设计：无需外部电源，功耗低。

  ● 短距离传输：通常支持1米到7米的传输距离。

  ● 高带宽：支持10Gbps、25Gbps、40Gbps、100Gbps甚至400Gbps的传输速率。

应用：
一个主要应用是通过SFP+接口连接网络硬件（数据中心机架内设备互联如服务器到交换机）。这种
连接能够在5米距离内以10Gbps的全双工速度传输数据。此外，这种配置的收发延迟比当前的
10GBASE-T Cat 6/Cat 6A/Cat 7电缆系统低15到25倍：SFP+双轴电缆的延迟为0.1微秒，而
10GBASE-T的延迟为1.5到2.5微秒。SFP+双轴电缆的功耗约为0.1瓦，远低于10GBASE-T的4到8
瓦。
有源光缆（AOC, Active Optical Cable）
有源光缆（AOC）是一种集成光模块的光纤电缆，将光模块的功能直接嵌入电缆两端，提供即插
即用的高速数据传输解决方案。
结构：
 ● 光纤：用于传输光信号。
 ● 光模块：集成在电缆两端，实现电光转换。

 ● 连接器：通常采用高速接口（如QSFP+、QSFP28）   。
 ● AOC的优势：

    ○ 高带宽：支持10Gbps到400Gbps的传输速率。
    ○ 低延迟：光信号传输速度快，延迟极低。

    ○ 轻便灵活：比铜缆更轻，适合高密度部署。

    ○ 即插即用：无需额外配置，安装简便。

 ● AOC的应用场景：
    ○ 数据中心内部互联（如服务器到交换机）     。
    ○ 高性能计算（HPC）和人工智能（AI）集群。

    ○ 短距离高速数据传输（通常≤100米）   。

AOC与DAC的对比
100G DAC & AOC Data Center Cabling Solutions | QSFPTEK
AOC和直接连接铜缆（DAC）都是数据中心常用的高速连接方案，但各有特点：
 特性      AOC（有源光缆）           DAC（直接连接铜缆）
 传输介质    光纤                  铜缆
 传输距离    较长（通常≤100米）         较短（通常≤7米）
 带宽      高（支持10Gbps到400Gbps） 高（支持10Gbps到400Gbps）
 抗干扰能力 强（不受电磁干扰）             较弱（受电磁干扰影响）
 成本      较高                  较低
 应用场景    数据中心内部互联、高性能计算 机架内设备互联
在数据通信和网络设备中，除了高速接口（如SFP、QSFP）之外，还有许多其他常见的接口类
型，例如 RJ45、USB、HDMI 等。这些接口在不同的应用场景中发挥着重要作用。

RJ45 是一种常见的以太网接口，采用8针模块化连接器设计，主要用于连接双绞线电缆（如
Cat5e、Cat6）。它支持从10/100Mbps快速以太网到10Gbps高速以太网的多种速率，传输距离通
常不超过100米。RJ45接口广泛应用于局域网（LAN）、家庭和企业网络以及网络设备（如交换
机、路由器和计算机）的连接，因其成本低、易于安装和广泛兼容性，成为以太网通信中最常用的
接口之一。
RS-232接口有22根线，采用标准的25芯插头座（DB-25）。除此之处，目前广泛应用的还有一种
9芯的RS-232接口（DB-9）。它们的外观都是一个D形的，不过对接的两个接口又分为针式的“公
头”和孔式的“母头”两种，DB-9“母头”和“公头”和DB-25的“母头”和“公头”分别如下图所示
USB（Universal Serial Bus） 是一种通用的串行总线接口，支持数据传输和设备供电。USB接口
有多种版本，包括USB 2.0（480Mbps）、USB 3.0（5Gbps）、USB 3.1（10Gbps）和USB 4
（40Gbps），广泛应用于计算机、外设（如键盘、鼠标、打印机）和移动设备（如智能手机、平
板）的连接，具有即插即用和热插拔的特点。
HDMI（High-Definition Multimedia Interface） 是一种高清多媒体接口，用于传输未压缩的视
频和音频信号。HDMI支持高达48Gbps的传输速率（HDMI 2.1），能够传输4K甚至8K分辨率的内
容，广泛应用于电视、显示器、游戏机和家庭影院系统，提供高质量的视听体验。
光纤接口（如LC、SC） 主要用于光纤通信，支持高带宽和长距离传输。LC接口体积小，适合高
密度部署；SC接口结构稳固，适合数据中心和电信网络。光纤接口广泛应用于数据中心互联、电
信骨干网和高性能计算环境。
无线媒体 or 非导向媒体
美国和中国的频谱划分在管理机制、分配方式和重点频段上存在显著差异。美国通过市场化拍卖和
动态共享机制，推动高频段（如毫米波）的开发；而中国则由国家主导，优先分配中低频段用于
5G建设，并支持行业专网发展。两国的频谱政策都旨在为5G和未来通信技术提供坚实的基础，但
侧重点和实施路径不同。




美国频谱划分                            中国频谱划分
美中频谱划分的对比
 方面     美国                           中国
 管理机构   联邦通信委员会（FCC）                 工业和信息化部（MIIT）
 分配机制   市场化拍卖                        国家统一规划
 5G重点频段 3.5 GHz（CBRS）、28 GHz、39      3.5 GHz、2.6 GHz、4.9 GHz
        GHz
高频段开发  积极推动毫米波（28 GHz、39              逐步推进毫米波（24.75 GHz～27.5
       GHz）                           GHz）
主要频段及应 低频段（<1 GHz）：                   低频段（<1 GHz）：
用       ● 600 MHz：用于5G和广播电             ● 700 MHz：用于5G和广播电视。
          视。                           ● 800 MHz：用于公共安全通信和
        ● 700 MHz：用于公共安全通信               4G LTE。
          和4G LTE。
                                     中频段（1 GHz～6 GHz）：
              中频段（1 GHz～6 GHz）：       ● 2.6 GHz：用于4G LTE和5G。
               ● 2.4 GHz：用于Wi-Fi、蓝牙和  ● 3.5 GHz：用于5G主流频段。
                 物联网（IoT）。            ● 4.9 GHz：用于5G和行业专网。
               ● 3.5 GHz（CBRS）：用于5G和
                 私有网络。
               ● 5 GHz：用于Wi-Fi和5G。
                                     高频段（>6 GHz）：
                                      ● 24.75 GHz～27.5 GHz：用于5G
              高频段（>6 GHz）：              毫米波通信。
               ● 28 GHz和39 GHz：用于5G毫
                                      ● 60 GHz：用于短距离高速通信
                 米波通信。                  （如WiGig）。
               ● 60 GHz：用于WiGig（高速短
                 距离通信）。

按照频率特性的信号划分




1.   地波传播（f < 2 MHz）
地波传播或多或少总沿着地球表面轮廓传播，并且能到达相当远的地方，超出视平线之外。大约在
2MHz以下的频率可以达到这个效果。这个频率内的电磁波之所以会沿着地表曲线传播是由多种因
素造成的。其中之一是电磁波会使地表出现感应电流，其结果是靠近地面的波阵面速度变慢，致使
波阵面向下倾斜，因而会沿着地表曲线传播。另一个因素是衍射，它是一种自然现象，与电磁波遇
到障碍物时的反应有关。在这个频率范围内的电磁波会被大气层散射，以至于它们无法穿透高层大
气。 地波通信广为人知的例子就是调幅无线电广播（AM radio）。
 2. 天波传播 （f < 2 MHz~30MHz）

天波传播用于业余无线电、CB无线电以及国际电台广播，如BBC和美国之音。使用天波传播时，
从基地天线发射出去的信号被高层大气的电离层反射回地球。虽然电波看起来好像被电离层反射，
就如同电离层是坚硬的反射面一样，但实际上这种效果是由折射引起的。稍后我们再来描述折
射。
天波信号可以经过多次反弹传播，在电离层和地球表面之间穿梭几个来回。在这种传播模式下，在
距离发送器几千千米以外的地方也能收到信号。
 3. 视距传播（f > 30MHz）

在高于30MHz的频段上地波和天波传播模式都无法工作，只有视距通信可行。对于卫星通信来
说，高于30MHz的信号不会被电离层反射，因此信号能够在地球站和位于地球站上空且在其视平
线以内的卫星之间传送，对于地面通信，发送天线和接收天线必须在双方有效视距以内。在这里使
用“有效”这个词是因为微波会被大气层弯曲和折射。 在30MHz到1GHz通常称作射频raido区域，
适用于全向通信。1GHz到40GHz叫做微波频率，适用于点对点传输，也用于卫星通信。

\section{[新增-2026-03] 高速光互联技术前沿}

随着AI大模型参数规模从百亿级向万亿级演进，GPU集群内部互联带宽需求持续攀升。448G及更高速率的光互联技术成为支撑下一代AI基础设施的关键技术之一。本节介绍高速光互联的技术要求和前沿发展趋势。

\subsection{448G互联技术背景}

AI大模型训练对网络带宽的需求呈指数级增长：
\begin{itemize}
    \item GPT-3级别模型：百Gbps级互联即可满足
    \item GPT-4级别模型：需要Tbps级互联能力
    \item 未来万亿参数模型：单端口需要448Gbps及以上速率
\end{itemize}

448G互联（即每通道56Gbps PAM4调制 $\times$ 8通道）是下一代高速互联的重要里程碑。

\subsection{448G互联技术要求}

\subsubsection{物理层要求}

\begin{table}[htbp]
\centering
\caption{448G光互联技术指标}
\begin{tabular}{|l|c|c|}
\hline
\textbf{指标} & \textbf{规格} & \textbf{说明} \\
\hline
单通道速率 & 56 Gbps & PAM4调制 \\
\hline
通道数量 & 8通道 & 并行传输 \\
\hline
总带宽 & 448 Gbps & 单向 \\
\hline
调制格式 & PAM4 & 每符号2比特 \\
\hline
信号速率 & 28 GBaud & 波特率 \\
\hline
传输距离 & 100m-2km & 根据应用场景 \\
\hline
误码率 & $<10^{-12}$ & 前向纠错后 \\
\hline
\end{tabular}
\end{table}

\subsubsection{光模块形态}

\textbf{QSFP-DD800}：
\begin{itemize}
    \item 支持8通道 $\times$ 100Gbps PAM4，总带宽800Gbps
    \item 向下兼容QSFP-DD和QSFP28
    \item 功耗控制在15W以内
\end{itemize}

\textbf{OSFP（Octal Small Form-factor Pluggable）}：
\begin{enumerate}
    \item 支持8通道，单通道速率可达112Gbps
    \item 支持液冷散热，适用于高功耗场景
    \item 成为800G/1.6T时代的主流形态
\end{enumerate}

\subsection{关键技术挑战}

\subsubsection{信号完整性挑战}

\begin{enumerate}
    \item \textbf{高频损耗}：56Gbps PAM4信号在PCB上损耗巨大，需采用低损耗材料
    \item \textbf{串扰控制}：高密度引脚间的信号串扰需精确控制
    \item \textbf{阻抗匹配}：全链路阻抗一致性要求极高
\end{enumerate}

\subsubsection{光电器件挑战}

\begin{itemize}
    \item \textbf{激光器带宽}：需要支持28GBaud及以上速率的EML或硅光调制器
    \item \textbf{探测器灵敏度}：高带宽下保持高灵敏度，降低功耗
    \item \textbf{封装散热}：高速器件功耗密度增加，需要先进散热方案
\end{itemize}

\subsubsection{前向纠错（FEC）}

448G互联需要更强的FEC能力：
\begin{enumerate}
    \item 采用级联FEC（CFEC）或拼接FEC（KP4+RS）
    \item 开销控制在10-15\%
    \item 支持AI/ML辅助的信道估计和自适应均衡
\end{enumerate}

\subsection{高速互联架构演进}

\subsubsection{从CPO到NPO}

\textbf{CPO（Co-Packaged Optics，共封装光学）}：
\begin{itemize}
    \item 将光引擎与交换芯片封装在一起
    \item 消除电信号在PCB上的长距离传输
    \item 降低功耗30\%以上，延迟减少50\%
\end{itemize}

\textbf{NPO（Near-Package Optics，近封装光学）}：
\begin{enumerate}
    \item 光引擎靠近交换芯片但独立封装
    \item 平衡性能和可维护性
    \item 有望成为中短期主流方案
\end{enumerate}

\subsubsection{全光交换架构}

\begin{itemize}
    \item 采用OXC实现端口间全光互联
    \item 消除光电转换环节，降低功耗和延迟
    \item 支持灵活的拓扑重配置
\end{itemize}

\subsection{智能高速接口}

\subsubsection{自适应信号处理}

\begin{enumerate}
    \item \textbf{自适应均衡}：根据信道特性自动调整均衡参数
    \item \textbf{动态功耗管理}：根据流量负载调节发射功率
    \item \textbf{温度补偿}：自动补偿温度变化引起的性能漂移
\end{enumerate}

\subsubsection{数字信号处理（DSP）}

\begin{itemize}
    \item 采用先进DSP算法补偿信道损伤
    \item 支持ML-based信道估计和均衡
    \item 实现更长的传输距离和更高的链路裕量
\end{itemize}

\subsubsection{链路诊断与遥测}

\begin{enumerate}
    \item 实时监控链路质量指标（BER、SNR、眼图裕量）
    \item 预测性维护，提前识别潜在故障
    \item 遥测数据用于网络级流量优化
\end{enumerate}

\subsection{技术演进路线图}

\begin{table}[htbp]
\centering
\caption{高速光互联技术演进}
\begin{tabular}{|l|c|c|c|c|}
\hline
\textbf{年代} & \textbf{速率} & \textbf{技术} & \textbf{调制格式} & \textbf{应用场景} \\
\hline
2023-2024 & 400G & QSFP-DD/QSFP112 & PAM4 56Gbps & 当前主流 \\
\hline
2025-2026 & 800G & QSFP-DD800/OSFP & PAM4 112Gbps & 大规模部署 \\
\hline
2027-2028 & 1.6T & CPO/OSFP-XD & PAM4 224Gbps & 下一代超算 \\
\hline
2030+ & 3.2T & 光电融合 & 高阶调制 & 未来演进 \\
\hline
\end{tabular}
\end{table}

高速光互联技术的持续演进为AI大模型训练提供了强大的带宽支撑，CPO/NPO、智能接口等创新技术将进一步提升互联效率和能源利用效率，推动AI基础设施向更高性能、更低功耗的方向发展。

\section{光通信技术：让光为网络插上翅膀}

人类通信的历史，是一部不断追求更快、更远、更大容量信息传递的历史。从古老的狼烟信号到电报、电话，再到现代的光纤通信，每一次技术的飞跃都深刻改变了信息社会的面貌。在众多传输介质中，光——这种我们日常生活中最熟悉的存在——却成为了当今信息时代最关键的传输载体。当我们浏览网页、观看视频、进行视频会议时，你是否想过，这些海量的数据是如何跨越千山万水，精准无误地送达你的屏幕？答案就是光通信技术。光纤通信作为20世纪最伟大的技术发明之一，已经成为现代通信网络的基石。从第一条20dB/km光纤的诞生到如今400G、800G高速光网络的规模部署，光通信技术用半个世纪的时间，重新定义了人类传递信息的方式。在本章中，我们将深入探讨光通信技术的发展脉络、核心技术原理以及未来演进方向，带你理解为什么说“光”才是网络的真正翅膀。

\subsection{光纤通信的诞生：从黑暗走向光明}

要理解光纤通信为什么能够取代传统的铜缆通信，我们首先需要回到20世纪60年代那个充满技术创新的年代。那时的电话网络主要依赖铜缆进行信号传输，随着通话需求的爆炸式增长，工程师们面临着一个严峻的物理瓶颈——铜缆的带宽极限。

让我们先思考一个问题：为什么电信号在金属导线中传输会有容量限制？这要从电磁学的基本原理说起。当高频电信号通过铜缆传输时，会遇到两个主要的物理限制。第一个是趋肤效应，随着信号频率的升高，电流趋向于在导体表面流动，导致有效导电面积减小，电阻增大，信号衰减加剧。第二个是串扰，即相邻导线之间的电磁耦合会引入噪声，限制可传输的信号质量。想象一下，在一根铜缆中同时传输成千上万路电话信号，它们相互干扰，最终能够清晰传递的距离和信号质量都受到严重制约。

在这样的背景下，科学家们开始思考：能否利用光来传输信息？光具有极高的频率——可见光的频率在430THz到750THz之间，这意味着理论上单根光纤可以携带极其庞大的信息量。相比之下，传统无线电波的频率仅为几十GHz到几百GHz，铜缆中的电信号频率更低，差距达到数万倍。然而，把这个理论潜力变为现实，科学家们花了将近十年时间，跨越了无数技术障碍。

1966年，华裔科学家高锟博士发表了一篇具有里程碑意义的论文《光频率介质纤维表面波导》，从理论上论证了利用光纤进行光通信的可能性。高锟指出，制约光通信发展的主要障碍不是光纤本身的损耗，而是当时玻璃材料的纯度不够。他认为，通过改进制造工艺，可以将石英光纤的损耗降低到20dB/km以下，从而实现远距离光通信。这篇论文在当时引起了巨大争议，因为许多科学家认为玻璃的损耗是固有的物理特性，无法克服。但高锟坚持自己的判断，并积极与玻璃制造商合作，推动实际工艺的改进。

1970年，一个值得铭记的年份。康宁玻璃公司（Corning Glass Works）宣布研制成功世界上第一根低损耗光纤，损耗率约为20dB/km，虽然这个数值距离理论极限还有很大距离，但已经足以让光通信成为可能。这一突破与同时期半导体激光器的发明一起，标志着光纤通信时代的正式开启。此后数十年间，光纤的损耗率不断降低，如今商用光纤的损耗约为0.2dB/km，这意味着光信号可以传输超过100公里才衰减到原来的一百分之一。

让我们用一个具体的数字来感受光纤的优越性。一根直径仅125微米（与头发丝相当）的单模光纤，在1550nm波长窗口的损耗为0.2dB/km。假设我们发射一个功率为1毫瓦（0dBm）的光信号，经过100公里传输后，接收端的功率仍有：
\[
P_{out} = P_{in} - \alpha \times L = 0\text{dBm} - 0.2\text{dB/km} \times 100\text{km} = -20\text{dBm}
\]
这相当于10微瓦的功率，在现代光电探测器灵敏度的范围内可以轻松检测到。换成同等传输条件的铜缆，信号早已被噪声淹没，无从辨识。

光纤通信之所以能够取得巨大成功，源于光信号独特的物理优势。首先是极低的损耗，我们已经看到，100公里仅20dB的总衰减意味着光信号可以传输极远的距离，这在跨城、跨省乃至跨洋通信中意味着可以大幅减少中继器的数量，降低系统成本和复杂度。其次是极宽的带宽，单模光纤的理论可用带宽超过10THz，相当于可以同时传输数亿路电话或数百万路高清视频。第三是抗电磁干扰能力，光信号在光纤中传输完全不受外部电磁场影响，因为光是一种电磁波，但它被束缚在光纤内部，不与外界发生电磁耦合。第四是保密性好，传统的铜缆传输会产生电磁辐射，容易被窃听，而光纤中的光信号被完全束缚在纤芯内部，除非物理破坏光纤，否则无法被窃听。最后是体积小、重量轻，一根光缆可以包含数百根光纤，而直径和重量远小于同等容量的铜缆，特别适合空间受限的管道铺设。

正是这些显著优势，使光纤通信在短短半个世纪内从实验室走向全球，从骨干网渗透到接入网，从电信运营商扩展到互联网公司，从传统通信领域进入数据中心。现在，光纤网络已经承载了全球超过90\%的通信流量，成为信息社会名副其实的“神经中枢”。

\subsection{光传输技术的演进：从PDH到OTN}

光纤的发明为光通信提供了物理基础，但要真正实现大容量、远距离的光信号传输，还需要一套完整的光传输技术体系。从上世纪70年代至今，光传输技术经历了多次重大变革，每一代技术都在解决前一代面临的根本性问题。理解这个演进过程，对于把握现代光网络的设计原理至关重要。

让我们首先理解一个核心概念：时分复用（TDM）。在数字通信中，我们要把大量独立的信号合并成一个数据流进行传输，最直观的方法是让这些信号“排队”使用传输通道——每个信号依次占用一小段时间，然后让下一个信号使用。想象一下，一根电话线同时有多人通话，系统会为每个人分配一个固定的时间片，在这个时间片内只有这个人的声音可以通过，时间片结束后切换到下一个人。这种“交替使用”的方式，就是时分复用的本质。在电通信时代，TDM技术被广泛用于提高铜缆的利用率。

准同步数字体系（PDH，Plesiochronous Digital Hierarchy）是最早的光传输技术之一。PDH技术出现于1970年代，它在光纤上传输数字信号，通过时分复用将多个低速信号合并成高速信号流。PDH的基本速率是2.048Mbps（欧洲标准）或1.544Mbps（北美标准），然后通过复用形成更高速率。典型的PDH复用结构是4:1复用，即4个2.048Mbps信号复用成一个8.448Mbps信号，再4个8.448Mbps复用成34.368Mbps，以此类推。

然而，PDH技术存在一个根本性的设计缺陷——时钟同步问题。由于历史原因，不同地区的电信运营商使用不同精度的时钟源，这些时钟虽然标称频率相同，但存在微小的偏差。为了正确分接信号，PDH需要在每个复用层级添加额外的开销比特来补偿时钟差异，这就导致了系统设计的复杂性。更麻烦的是，当我们需要从高速信号中提取某一低速支路信号时，必须将高速信号完全解复用，这种“解复用-再复用”的过程会引入抖动和漂移，影响信号质量。

为了克服PDH的这些缺点，同步数字体系（SDH，Synchronous Digital Hierarchy）应运而生。SDH最初由美国贝尔通信研究所（Bellcore）提出，后来被ITU-T正式标准化为同步数字体系（SDH）和同步光纤网络（SONET）。SDH的核心创新在于引入了同步机制——网络中所有设备都使用同一个高精度时钟源，确保各节点严格同步。这样，SDH可以直接从高速信号中分接出任意低速支路，无需像PDH那样进行繁琐的逐级解复用。

SDH的帧结构是其技术精髓所在。SDH帧是一个结构化的信息块，以STM-1（155.52Mbps）为基本模块，更高速率则由多个STM-1通过字节间插同步复用而成。STM-1帧结构为9行×270列字节，传输周期为125微秒，即每秒传输8000帧，这与PCM语音信号的采样率完美匹配。帧结构中包含了丰富的开销字节，用于网络管理、性能监控、故障定位等功能。特别是段开销（Section Overhead）和通道开销（Path Overhead）的设计，使SDH网络具备了强大的OAM（操作、管理、维护）能力。

SDH技术在1990年代得到了广泛应用，我国的主要运营商都建设了基于SDH的骨干传输网。SDH的环形网络拓扑和自动保护倒换（APS）机制，使其具备了极高的可靠性——当光环网上某处发生故障时，业务可以在50毫秒内自动切换到备用路径，满足了语音通信等实时业务对网络可靠性的严格要求。

然而，随着互联网的爆发式增长，SDH的局限性逐渐显现。SDH是针对传统语音业务设计的，它的带宽分配是刚性的——一旦为某个业务分配了STM-1的带宽，无论该业务是否实际使用这部分带宽，其他业务都无法利用。这种“静态管道”的模式在以IP为代表的突发性数据业务面前显得效率低下。同时，SDH对分组业务的支持能力有限，虽然后来出现了MSTP（Multi-Service Transport Platform，多业务传输平台）来在SDH上承载以太网和IP业务，但这只是权宜之计。

波分复用（WDM，Wavelength Division Multiplexing）技术的出现，为光传输开辟了全新的方向。WDM的原理与无线电通信中的频分复用类似——将不同波长的光信号合并到同一根光纤中传输，在接收端再将它们分开。由于不同波长的光信号彼此独立，可以把它们看作独立的“光通道”，每条通道承载独立的业务。这种“空间换频率”的策略，使单根光纤的容量得到了成倍的提升。

粗波分复用（CWDM，Coarse WDM）是最早的WDM技术，它使用间隔较宽（20nm）的波长通道，典型配置是18个波长通道。由于波长间隔宽，CWDM对激光器的波长精度要求较低，成本相对低廉，但容量有限，主要用于城域网或接入网。

密集波分复用（DWDM，Dense WDM）则将波长间隔压缩到0.4nm甚至更窄，在C波段（1525nm-1565nm）可以同时传输80-160个波长通道，每个波长通道可以承载10Gbps、40Gbps或100Gbps的信号，使单根光纤的总容量达到数Tbps甚至数十Tbps。DWDM技术彻底改变了长距离传输的经济性——原本需要数十条光纤才能承载的流量，现在一条光纤就可以搞定。

然而，DWDM网络也面临新的挑战。在没有电中继的情况下，DWDM信号可以传输的距离受限于光纤损耗和非线性效应。对于100Gbps以上的高速信号，传输距离进一步缩短，需要在每隔80-100公里设置光放大器（如EDFA，Erbium-Doped Fiber Amplifier，掺铒光纤放大器）和色散补偿装置。更重要的是，DWDM系统缺乏像SDH那样的标准化开销机制，网络管理、故障定位、带宽调度等能力都较弱。

光传送网（OTN，Optical Transport Network）正是为了解决这些问题而诞生的。OTN由ITU-T G.709标准定义，它在DWDM的光层之上引入了一个独立的电层处理机制，为波长通道提供了类似SDH的帧结构和开销管理能力。OTN帧的结构称为ODU（Optical Data Unit），从ODU1到ODU4，对应不同的比特速率。ODU4的速率约为100Gbps，与100G DWDM波长通道相匹配。

OTN的核心价值在于提供了一种面向未来的传输架构。它继承了SDH的强大OAM能力——G.709帧结构中包含了丰富的开销字段，用于性能监控、告警指示、故障定位等；同时，它支持波长级的大容量传输，能够充分发挥DWDM技术的带宽潜力。此外，OTN还支持多种客户信号映射（如以太网、SDH、FC等），是一种真正的多业务承载平台。

在OTN架构中，电交叉和光交叉是两种主要的交换机制。电交叉发生在ODU层面，通过电信号处理实现客户信号的复用、解复用和调度，典型设备是OTN交换机或电交叉芯片。光交叉（ROADM，Reconfigurable Optical Add-Drop Multiplexer）则在光层实现波长级别的路由选择——通过波长选择开关（WSS）技术，可以在任意波长端口上下载或穿通任意波长的信号，实现了波长级别的灵活调度。ROADM的出现，使光网络具备了前所未有的灵活性，被认为是下一代智能光网络的核心设备。

近年来，400G DWDM技术开始规模部署。相比100G，400G采用更高级的调制格式（如16QAM或64QAM），在相同的频谱资源下将单波长容量提升了4倍。但高级调制格式对信噪比要求更高，传输距离相应缩短，因此400G系统主要应用于城域和数据中心互联场景。800G和1.6T的研发也已在积极推进中。

\begin{table}[htbp]
\centering
\caption{光传输技术演进对比}
\begin{tabular}{|l|c|c|c|c|}
\hline
\textbf{技术体制} & \textbf{典型速率} & \textbf{复用方式} & \textbf{主要优势} & \textbf{典型应用} \\
\hline
PDH & 2-140Mbps & 时分复用 & 早期标准，简单 & 已被淘汰 \\
\hline
SDH & 155Mbps-10Gbps & 同步时分复用 & 高可靠性，强OAM & 传统骨干网 \\
\hline
WDM/CWDM & 10Gbps×n & 波分复用 & 容量大 & 城域网 \\
\hline
DWDM & 10/40/100Gbps×80+ & 密集波分复用 & 超大容量 & 骨干网/海缆 \\
\hline
OTN & 10/100/400Gbps & 光电复合 & 弹性带宽，智能调度 & 下一代网络 \\
\hline
\end{tabular}
\end{table}

从上表可以看出，光传输技术的演进主线是不断提高单波长速率、增加波长通道数、提升网络灵活性。每一代技术都有其特定的历史使命和应用场景，PDH和SDH为光纤通信的早期发展奠定了基础，WDM/DWDM实现了传输容量的飞跃，OTN则代表了光网络智能化的方向。

\subsection{光纤接入网：从“最后一公里”到全光网络}

前面我们讨论的光传输技术，主要应用在骨干网和城域网层面，承担着城市之间、区域之间的远距离信息传递任务。然而，对于普通用户来说，无论是用手机刷视频还是在家上网，数据的起点和终点往往在用户终端——这就涉及到了接入网的概念。

接入网是通信网络中离用户最近的部分，它负责将用户终端连接到运营商的核心网络。由于直接面向海量用户，接入网的技术选择需要兼顾带宽、覆盖、成本和运维等多个维度。在众多接入技术中，基于光纤的PON（Passive Optical Network，无源光网络）技术已经成为当前主流的宽带接入方案。

让我们首先理解一个基本概念：为什么叫“无源”光网络？传统的电话网络或早期宽带网络，在从运营商机房到用户家庭的这段距离中，往往需要部署大量的有源设备——如放大器、交换机等，这些设备需要供电、维护，成本较高。而PON网络采用无源分光器来实现光信号的分配，这些分光器不需要电源、不需要维护，从而大幅降低了网络建设和运维成本。

一个典型的PON系统由三部分组成：光线路终端（OLT，Optical Line Terminal）位于运营商机房，光网络单元（ONU，Optical Network Unit）位于用户侧，光分配网络（ODN，Optical Distribution Network）连接OLT和ONU。OLT是整个PON系统的核心，负责向上连接城域骨干网，向下管理所有的ONU设备。ONU则安装在用户家中或楼道，提供以太网、Wi-Fi等用户接口。

在典型的树形拓扑中，从OLT发出的光信号经过1:N的分光器后，分发给N个ONU。这里的N通常是8、16、32或64，取决于分光比和传输距离。关键在于，下行方向（从OLT到ONU）的广播式传输——所有ONU都能收到所有数据，但每个ONU只提取属于自己的部分；上行方向（从ONU到OLT）则采用时分多址（TDMA）技术，各个ONU在不同的时间槽中发送数据，避免冲突。

 PON技术经历了从APON、BPON到GPON和EPON的演进。APON（ATM PON）是最早的PON标准，基于ATM协议，但由于ATM协议的复杂性，最终没有获得大规模应用。BPON（Broadband PON）在APON基础上增加了动态带宽分配等特性，但同样未能普及。

GPON（Gigabit PON）和EPON（Ethernet PON）是目前主流的两种PON技术，分别由两个不同的标准化组织主导。GPON由ITU-T G.984系列标准定义，属于国际电信联盟体系，下行速率最高2.5Gbps，上行最高1.25Gbps，支持非对称速率配置。GPON的效率较高——由于采用了GEM（G-PON Encapsulation Mode）封装方式，可以更高效地承载以太网帧和TDM业务。EPON由IEEE 802.3ah标准定义，属于电气电子工程师协会体系，采用以太网帧结构，与现有IP网络天然兼容，部署成本相对较低。这两大标准体系在技术路线上各有侧重：GPON在带宽效率、QoS保障和覆盖距离方面更有优势，而EPON在成本控制和与以太网兼容性方面更具竞争力。我国运营商根据自身网络条件和业务需求，选择了不同的技术路线——中国电信和中国联通主要部署GPON，中国移动则大规模推广EPON，近年来又都向10G-PON演进。

\begin{table}[htbp]
\centering
\caption{主流PON技术对比}
\begin{tabular}{|l|c|c|c|c|}
\hline
\textbf{技术标准} & \textbf{下行速率} & \textbf{上行速率} & \textbf{分光比} & \textbf{主要优势} \\
\hline
EPON & 1Gbps & 1Gbps & 1:32 & 成本低，兼容以太网 \\
\hline
GPON & 2.5Gbps & 1.25Gbps & 1:64 & 效率高，覆盖远 \\
\hline
10G-EPON & 10Gbps & 10Gbps & 1:64 & 带宽提升10倍 \\
\hline
XG(S)-PON & 10Gbps & 2.5Gbps/10Gbps & 1:128 & 超高带宽，下一代主流 \\
\hline
50G-PON & 50Gbps & 25/50Gbps & 1:256 & 面向未来超宽带 \\
\hline
\end{tabular}
\end{table}

10G-PON是PON技术的重大飞跃，它将下行速率提升到10Gbps，是GPON的4倍、EPON的10倍。10G-PON包括XG-PON（非对称，10G/2.5G）和XGS-PON（对称，10G/10G）两种版本。目前，10G-PON正在全球范围内规模部署，我国运营商的用户覆盖已经超过1亿户。

展望未来，50G-PON已被ITU-T标准化，预计将在2025-2030年逐步商用。50G-PON将提供50Gbps的下行速率，足以满足8K视频、AR/VR、云游戏等新兴应用的需求。更重要的是，50G-PON实现了与现有GPON/XGS-PON的共存——可以在同一ODN网络上逐步升级，保护运营商的既有投资。

FTTH（Fiber To The Home，光纤到户）是接入网的终极目标——将光纤直接敷设到用户家庭，用户独享全部带宽资源。我国是全球FTTH部署最领先的国家，根据工信部数据，截至2024年底，我国光纤宽带用户占比已超过95\%，百兆以上宽带用户占比超过90\%，千兆宽带用户数突破2亿。

从技术角度看，FTTH的实现方式包括单星型拓扑和树形拓扑两种。单星型拓扑中，每个用户从OLT直接拉一条专属光纤到户，带宽独享、延迟最低，但光纤资源消耗大，适合高价值商业用户。树形拓扑采用分光器共享光纤资源，建设成本低，是目前主流的FTTH部署模式。

从“最后一公里”到全光网络，光纤接入技术一直在持续演进。回顾这段历史，我们可以看到一个清晰的趋势：带宽越来越宽、距离越来越远、成本越来越低、用户体验越来越好。PON技术以其独特的设计理念——用无源器件替代有源设备——创造了一种极具性价比的宽带接入方案，正在并将持续改变亿万用户的网络体验。值得一提的是，我国是全球FTTH部署最领先的国家，根据工信部数据，截至2024年底，我国光纤宽带用户占比已超过95\%，百兆以上宽带用户占比超过90\%，千兆宽带用户数突破2亿，遥遥领先于其他国家和地区。

\subsection{数据中心光互联：高速网络的核心支撑}

如果说骨干网是通信的“高速公路”，那么数据中心内部网络就是数据的“城市交通网”。随着云计算、大数据和人工智能的快速发展，数据中心已成为现代信息基础设施的核心据点。在这些庞大的计算集群中，成千上万的服务器和存储设备需要高速互联，形成一个高效的信息交换平面。而光互联技术，正是支撑这个平面的关键技术。

让我们首先理解数据中心网络的连接需求。一个典型的数据中心包含数以万计的服务器，它们分布在多个机架（rack）中。每个机架顶部有一台ToR（Top of Rack）交换机，连接机架内所有服务器；ToR交换机再通过上行链路连接到汇聚层交换机；汇聚层交换机进一步连接到核心交换机。这种三级网络架构（接入-汇聚-核心）构成了数据中心网络的经典拓扑。在典型的带宽分配中，接入层（ToR）提供服务器上联带宽，汇聚层承上启下，核心层负责南北向出口，三层的带宽配比通常遵循40\%:20\%:40\%的原则——即ToR上行和核心下行的带宽各占40\%，汇聚层的南北向流量占20\%，这种设计确保各层都不会成为流量瓶颈。

在这些连接中，光模块扮演着至关重要的角色。光模块是实现电信号与光信号相互转换的可插拔器件，它通常包括发射器（激光器+驱动电路）、接收器（探测器+放大电路）和数字诊断功能。根据传输距离和应用场景的不同，光模块可分为多种类型。

短距离光模块主要用于机架内部或相邻机架间的连接，典型传输距离在100米以内。这类场景通常采用多模光纤，850nm波长的垂直腔面发射激光器（VCSEL）是主要的光源方案。多模光纤的优势在于对接精度要求较低，连接器成本低廉，非常适合数据中心内部短距离传输。典型的产品如SFP+（10G）、SFP28（25G）、QSFP28（100G）等。

长距离光模块用于数据中心之间或机架间较长距离的连接，典型传输距离在500米到80公里之间。这类场景主要采用单模光纤，1310nm或1550nm波长的分布反馈（DFB）激光器是主要光源。单模光纤的纤芯直径仅9微米，对接精度要求高，但传输损耗低、色散小，适合远距离传输。

在数据中心网络中，不同距离的连接需要采用不同的光模块技术。让我们以一个具体的例子来说明。假设一个数据中心有1000台服务器，每台服务器配备25G网卡，ToR交换机为48×25G下行+4×100G上行配置。那么服务器到ToR的距离通常在5-30米之间，采用25G AOC（有源光缆）或SFP28光模块；ToR到汇聚交换机的距离可能在100-300米之间，采用100G QSFP28光模块（多模SR4或单模CWDM4/PSM4）；如果汇聚交换机与核心交换机不在同一机房，距离可能达到数公里，就需要采用相干光模块或100G LR4（10公里单模）。

\begin{table}[htbp]
\centering
\caption{数据中心光模块技术对比}
\begin{tabular}{|l|c|c|c|c|c|}
\hline
\textbf{封装形式} & \textbf{速率} & \textbf{光纤类型} & \textbf{典型距离} & \textbf{应用场景} & \textbf{成本} \\
\hline
SFP+ & 10G & 多模/单模 & 300m/10km & 服务器接入 & 低 \\
\hline
SFP28 & 25G & 多模/单模 & 100m/10km & 服务器接入 & 中 \\
\hline
QSFP28 & 100G & 多模SR4 & 100m & ToR互联 & 中 \\
\hline
QSFP28 & 100G & 单模CWDM4 & 2km & 汇聚层 & 中 \\
\hline
QSFP28 & 100G & 单模PSM4 & 500m & 短距离单模 & 中 \\
\hline
QSFP-DD & 400G & 多模/单模 & 100m-10km & 交换机互联 & 高 \\
\hline
OSFP & 400G/800G & 单模 & 2-10km & 核心网络 & 高 \\
\hline
\end{tabular}
\end{table}

400G光模块是当前数据中心网络的主流技术。2019年，400G光模块开始规模部署，主要形态包括QSFP-DD（Double Density）和OSFP。QSFP-DD保持了与QSFP28兼容的尺寸，但通过8×50G的电通道实现400G速率；OSFP（Octal Small Form Factor Pluggable）是新定义的封装标准，尺寸稍大，但支持8×50G通道并预留了更高的功率预算。

在调制技术方面，PAM4（Pulse Amplitude Modulation 4-level，四电平脉冲幅度调制）已成为400G及以上速率的核心技术。相比传统的NRZ（Non-Return-to-Zero）调制，PAM4每个符号周期可以传输2比特信息，使有效数据速率翻倍。以50G PAM4为例，符号率为28Gbaud，但有效数据速率达到56Gbps（28G×2）。这种高效的调制方式使现有电接口和光纤基础设施可以承载更高的比特率。

2024-2025年，800G光模块开始大规模部署。800G采用8×100G PAM4的方案，每个通道使用112Gbaud的PAM4调制。在形态上，800G光模块主要有QSFP-DD800和OSFP两种标准。800G光模块主要用于AI训练集群的交换机互联，以及超大规模数据中心的东西向流量承载。

展望未来，部分厂商预计1.6T光模块将在2026-2027年开始商用。1.6T将采用8×200G或16×100G的方案，更高速率版本可能采用200G波特率的PAM4或更高阶的调制格式。同时，CPO（Co-Packaged Optics，共封装光学）技术正在成为研发热点——将光引擎与交换芯片封装在一起，取代传统的可插拔光模块，可以显著降低功耗和延迟。

对于AI训练集群来说，光互联的重要性更加突出。现代AI训练通常使用数百甚至数千块GPU，这些GPU之间需要频繁交换梯度数据和模型参数。以一个典型的 trillion参数大模型训练任务为例，单次迭代可能产生TB级的通信流量。如果网络延迟过高或带宽不足，GPU将处于等待通信完成的空闲状态，严重影响训练效率。因此，AI训练集群对网络的要求不仅是高带宽，更是超低延迟和确定性时延。

NVLink和InfiniBand已成为AI训练网络的主流技术。NVLink是NVIDIA开发的高速GPU互连技术，通过私有协议实现GPU间的直接通信，延迟可低至微秒级。当NVLink扩展到多GPU服务器之间时，就需要借助InfiniBand网络或基于InfiniBand的NVLink over Fabric。2024年，NVIDIA推出了Spectrum-X以太网平台，专门针对AI训练场景优化，结合了RDMA和自适应路由等技术。

数据中心光互联技术的发展，折射出整个ICT行业的演进趋势：从10G、25G到100G、400G、800G乃至未来的1.6T，速率不断提升；从NRZ到PAM4，调制效率持续优化；从可插拔模块到CPO，封装形态不断创新。每一次技术突破，都推动着数据中心网络向更高性能、更低功耗的方向发展，为云计算、大数据和人工智能提供更强大的基础设施支撑。

\subsection{FIO-IP与OTN跨层互联：打破光与IP的边界}

在现代通信网络中，IP和OTN是两个不同层面的技术——IP相当于通信网络中的“邮政编码系统”，负责确定数据包从源头到目的地的路由路径；OTN则相当于“物流公司的运输网络”，负责实际搬运这些包裹的卡车、飞机和仓库。IP层面关注的是“你的数据要去哪里”，OTN层面关注的是“通过什么路线、用什么方式把数据运过去”。两者需要紧密配合，才能提供端到端的业务服务。这种跨层协同，正是网络架构设计中的一大挑战。

传统的网络设计中，IP层和OTN层相对独立——IP路由器通过高速光接口直连光纤或WDM设备，OTN网络作为透明管道传递IP流量。这种简单直接的连接方式在早期网络中运行良好，但随着网络规模扩大和业务需求复杂化，跨层网络的运维挑战日益凸显。

让我们通过一个具体例子来说明问题。假设某运营商骨干网有多个核心节点，每个节点配置了多台IP路由器和OTN设备，它们之间通过光纤形成复杂的Full Mesh连接。IP层面，这些路由器通过BGP协议形成全互联拓扑，任意两台路由器之间都有业务路径。OTN层面，每条IP链路可能映射到不同的OTN波长或ODU通道，不同的波长可能经过不同的物理路由。

在日常运维中，这样的跨层网络会遇到几个典型问题。第一是电路错连——由于光纤连接错误或标签缺失，本应连接到节点A的光纤被错误地连接到了节点B，导致IP层面的路由计算结果与实际物理连接不符。第二是单点故障——如果OTN层面的两条业务路径恰好使用了相同的光纤或同一个光放大器，一旦该设备故障，将导致IP层面两条独立路径同时中断，失去冗余保护的意义。第三是资源利用率不均——IP层面的负载可能是均衡的，但OTN层面的波长资源可能分布不均，导致部分波长拥塞而其他波长空闲。

为了解决这些问题，业界发展出了多种跨层协同技术。FIO（Fabric Intelligence Orchestrator）是阿里巴巴网络研发团队开发的一套IP-OTN跨层协同系统，它的核心理念是打通IP层和OTN层的信息孤岛，实现跨层拓扑的统一视图和协同优化。

FIO的核心功能包括三个方面。首先是跨层拓扑发现，FIO通过收集IP网络的IGP拓扑信息和OTN网络的GMPLS控制平面信息，构建完整的跨层网络拓扑。在这个拓扑中，每个IP链路可以追溯到其对应的OTN路径，每个OTN波长可以关联到其承载的IP业务。这种端到端的视图是跨层分析的基础。

其次是可靠性分析。在跨层拓扑的基础上，FIO可以分析IP层面的业务路径在OTN层面的物理依赖关系。系统会自动识别“殊途同归”的情况——即IP层面看似冗余的多条路径，在OTN层面实际上共享了相同的物理资源。一旦识别出这种风险，系统会生成整改建议，建议运营商调整OTN层面的波长分配或光缆路由，消除单点故障。

第三是自动化的资源调度。FIO可以根据IP层面的流量预测和业务需求，自动计算OTN层面的波长和时隙分配方案，实现跨层资源的协同优化。这种自动化能力大大减轻了运维人员的工作负担，提高了网络资源利用效率。

FlexE（Flexible Ethernet，灵活以太网）是另一种重要的跨层技术，由OIF（Optical Internetworking Forum）定义。FlexE在传统以太网和OTN之间架起了一座桥梁——它将以太网帧切片映射到OTN的时隙中，实现了以太网业务对OTN带宽的灵活共享。

FlexE的核心概念包括FlexE Group和FlexE Client。一个FlexE Group由多个物理以太网接口组成，这些接口在FlexE层面被捆绑成一个逻辑的高速管道。FlexE Client是映射到FlexE Group中的子速率客户信号，它可以占用Group中特定的时隙，实现带宽的灵活分配——某个FlexE Client可以独享一部分带宽，也可以与其他Client共享带宽。

FlexE的价值在于提供了一种介于以太网和OTN之间的中间层技术。对于需要OTN级别可靠性和端到端管理的大型租户专线业务，FlexE可以提供硬管道隔离，确保业务带宽不被其他流量抢占。同时，FlexE保持了以太网的兼容性，业务无需感知底层传输机制的复杂性。

在5G移动回传网络中，FlexE也发挥着重要作用。5G网络对带宽和时延的要求远高于4G，同时需要支持网络切片。FlexE可以提供基于时隙的硬隔离，为不同类型的5G业务（如eMBB、URLLC、mMTC）提供差异化的传输保障。

G.709 OTN帧结构是OTN技术的核心，它定义了如何将各种客户信号（如以太网、SDH、FC等）封装进OTN帧，以及如何通过丰富的开销字段实现网络管理。OTN帧的基本结构包括帧对齐开销、段开销、通道开销和净荷区域。帧对齐开销用于帧定界和同步；段开销（SOH）负责光传输段层的监控和管理；通道开销（POH）负责端到端的客户信号监控；净荷区域则承载实际的用户数据。

OTN的复用结构采用逐级复用方式：客户信号首先映射进OPU（Optical Payload Unit），OPU再封装进ODU（Optical Data Unit），ODU进一步封装进OTU（Optical Transport Unit）。这种层级化的结构使OTN可以灵活承载从1Gbps到400Gbps的各种业务，同时保持与现有SDH/SONET网络的兼容性。

从FIO到FlexE再到G.709，跨层互联技术的发展反映了网络架构演进的大趋势——打破传统分层架构的边界，构建更加融合、智能的端到端网络能力。在未来，随着网络规模继续扩大、业务需求持续复杂化，跨层协同技术将变得更加重要。

\subsection{光互连技术前沿：从可插拔到光电融合}

经过半个世纪的发展，光通信技术已经取得了巨大的进步。然而，随着数据中心的规模持续扩张、AI训练集群对带宽的需求急剧增长，传统可插拔光模块的局限性日益凸显。功耗、密度、成本——这些挑战推动着光互连技术向新的方向演进。

让我们首先理解可插拔光模块面临的困境。现代数据中心交换机正朝着51.2T乃至更高带宽演进。以51.2T交换机为例，假设采用128×400G的端口配置，需要128个400G光模块。这些光模块的功耗、散热和空间占用都成为系统设计的瓶颈。一个典型的400G QSFP-DD光模块功耗约为10-15W，128个模块的总功耗就超过1kW，加上交换机的功耗，整个系统的散热需求十分惊人。

可插拔光模块的另一个局限是电接口瓶颈。从交换芯片到光模块的电气通道需要驱动56Gbaud甚至112Gbaud的PAM4信号，这些高速信号的传输距离受限于PCB的走线长度和损耗。当光模块数量众多、面板密度很高时，PCB布局和信号完整性设计变得极其困难。

共封装光学（CPO，Co-Packaged Optics）技术正是为了解决这些问题而提出的。CPO的核心理念是将光引擎（Optical Engine）——包括激光器、调制器、探测器等光电器件——与交换芯片（Switch ASIC）集成在同一个封装基板上。这样，光信号在芯片内部产生后，直接通过封装基板上的高速通道传输到光纤接口，消除了传统设计中PCB走线的限制。

CPO技术的优势主要体现在三个方面。第一是降低功耗——由于消除了驱动PCB电容的功耗，CPO可以将光互连的能效提升30\%-50\%。第二是减小延迟——光信号无需经过PCB上的长距离传输，信号延迟可减少50\%以上。第三是提高密度——光引擎与芯片的直接集成，使得在相同面板空间内可以支持更多的光端口。

然而，CPO技术也面临重大挑战。首先是散热问题——交换芯片和光引擎集成在同一封装中，散热设计极其复杂。其次是维护性——CPO模块一旦故障，需要更换整个交换系统模组，不能像可插拔光模块那样热插拔快速修复。第三是供应链——CPO需要光电器件和芯片的深度整合，对供应链整合能力要求很高。

正因为CPO的这些挑战，近封装光学（NPO，Near-Package Optics）作为一种折中方案受到关注。NPO将光引擎放置在交换芯片附近（而不是同一封装中），通过高速柔性电路板（High-Speed Flex）连接。与CPO相比，NPO在散热和维护性方面更加可控，同时仍能获得显著的功耗和延迟改善。

硅光子学（Silicon Photonics）是推动CPO/NPO发展的底层技术。传统光模块使用磷化铟（InP）或砷化镓（GaAs）作为光电器件材料，这些材料与CMOS半导体工艺不兼容，无法实现大规模集成。硅光子学使用成熟的CMOS工艺在硅片上制造光器件，包括硅波导、调制器、探测器等。虽然硅的光电性能不如InP，但硅光子技术的最大优势是能够利用成熟的半导体制造工艺，实现光电器件与控制电路的单片集成，大幅降低成本和尺寸。

硅光子在数据中心光模块中已经开始应用。Intel、Broadcom等公司已推出基于硅光子的100G、400G光模块产品。硅光子的集成优势使其在高速率（400G及以上）场景更具竞争力，预计到2028年，硅光子光模块的市场份额将超过50\%。

光电路交换（OCS，Optical Circuit Switch）代表了另一种光网络创新方向。传统的电交换需要将光信号转换为电信号进行处理，消耗大量功耗和时延。OCS在光层实现端口间的直连切换，数据始终以光的形式传输，无需光电转换。OCS特别适合需要频繁改变网络拓扑的数据中心场景，可以显著提高网络灵活性和能效。

阿里云在光互连领域的实践值得关注。从2015年开始，阿里巴巴启动了光模块自研项目，逐步掌握了光模块的设计和制造能力。目前，阿里巴巴的数据中心已部署了数百万只自研光模块，覆盖了从10G到400G的全系列 产品。通过光模块自研，阿里巴巴不仅显著降低了采购成本，更重要的是摆脱了对单一设备厂商的依赖，掌握了技术演进的主导权。在400G时代，阿里主导成立了SFP-DD MSA和QSFP112 MSA标准组织，联合上下游厂商共同制定技术规范，推动了行业技术标准的统一。目前，阿里巴巴正在积极布局800G/1.6T光模块、硅光子技术和CPO共封装光学，目标是实现从光芯片到光模块的全链路自主可控。根据公开信息，阿里巴巴已累计出货超过100万只自研光模块，成为全球光模块用量最大的互联网公司之一。

展望未来，光互连技术的发展将呈现几个明确趋势。第一是带宽持续提升——从400G到800G再到1.6T，单端口速率将继续翻倍增长。第二是光电融合加深——从可插拔到CPO/NPO再到完全光电集成，集成度不断提高。第三是智能化——光网络将具备更强的状态感知、故障预测和自适应调整能力。第四是绿色化——随着“双碳”目标的推进，降低功耗将成为光互连技术的核心追求。

\begin{table}[htbp]
\centering
\caption{光互连技术演进路线}
\begin{tabular}{|l|c|c|c|c|}
\hline
\textbf{技术阶段} & \textbf{时间节点} & \textbf{典型速率} & \textbf{核心特点} & \textbf{应用场景} \\
\hline
可插拔模块 & 2015-2024 & 100G-400G & 标准化的模块化产品 & 当前主流部署 \\
\hline
QSFP-DD800 & 2024-2025 & 800G & PAM4 112G SerDes & AI训练集群 \\
\hline
OSFP-XD & 2026-2027 & 800G-1.6T & 更高功率预算 & 下一代数据中心 \\
\hline
NPO & 2025-2027 & 1.6T-3.2T & 光引擎靠近芯片 & 过渡方案 \\
\hline
CPO & 2027+ & 3.2T+ & 光电同封，极致能效 & 超大规模集群 \\
\hline
\end{tabular}
\end{table}

光通信技术之所以能不断发展创新，根本动力来自于人类对更高速率、更大容量、更低延迟通信的永恒追求。从第一根低损耗光纤诞生到如今，CMOS硅光子集成、智能光网络等创新技术正在重塑光通信的形态。可以预见，在未来相当长的时期内，光仍将是网络世界最重要的信息载体，光通信技术将继续沿着更高性能、更低功耗、更智能的方向演进。

\subsection{本章小结}

光通信技术是现代信息社会的基础设施，从光纤通信的诞生到如今多元化应用，已经走过了半个多世纪的历程。在本章中，我们系统地探讨了光通信技术的发展脉络、核心技术原理和未来演进方向。

首先，我们回顾了光纤通信的诞生历程。从1966年高锟博士提出光纤通信的理论可能性，到1970年康宁公司研制成功20dB/km低损耗光纤，光纤通信经历了从理论到实践的关键突破。光信号相比电信号具有损耗低、带宽宽、抗电磁干扰、保密性好等独特优势，使其成为当代通信网络的理想传输介质。

其次，我们梳理了光传输技术的演进路线。从PDH到SDH，从WDM到DWDM再到OTN，每一代技术都解决了前一代的核心痛点。PDH/SDH实现了光纤上的数字化传输，WDM/DWDM大幅提升了单根光纤的传输容量，OTN则在光层之上引入了类似SDH的帧结构和开销管理能力，实现了光网络的智能化。

在接入网领域，PON技术以其“无源”的独特设计理念，成为光纤到户的主流方案。从GPON/EPON到10G-PON再到未来的50G-PON，PON技术持续演进，为亿万用户提供了越来越宽带的网络接入体验。

数据中心光互联是近年来发展最快的领域之一。400G光模块已规模部署，800G光模块正在大规模商用，1.6T光模块和CPO技术也在积极推进中。AI训练等新兴应用对网络带宽和延迟提出了更高要求，推动着光互联技术不断突破极限。

跨层互联技术（FIO、FlexE、G.709 OTN）打通了IP层和OTN层的信息孤岛，实现了端到端的网络视图和协同优化。随着网络规模扩大和业务需求复杂化，跨层协同将成为网络架构设计的重要方向。

最后，光互连技术正朝着光电融合的方向演进。CPO/NPO技术将光引擎与交换芯片集成，可大幅降低功耗和延迟；硅光子技术为光电集成提供了CMOS兼容的工艺平台；光电路交换（OCS）实现了纯光层面的网络切换。这些前沿技术创新，将推动光通信进入新的发展阶段。

光通信技术不仅仅是通信工程师的专业领域，它已经渗透到我们日常生活的方方面面。当我们刷视频、逛网购、玩云游戏时，光纤正在以光速传递着每一个比特的信息。未来，随着AI、元宇宙、量子通信等新技术的发展，光通信将面临新的机遇与挑战。但无论如何演进，“让光为网络插上翅膀”这一核心理念将始终指引着光通信技术的发展方向。
